First suppose that \(h\in V_m\). By the definition of \(V_m\), there are real coefficients \((\gamma_c)_{c\in\mathcal I_m}\) such that
\[
h(q)=\sum_{c\in\mathcal I_m}\gamma_c B_c^{(m)}(q)
\qquad(q\in K_{\epsilon}).
\]
Let \(b\in\mathcal B_{m,\epsilon}\) be feasible and let \(\nu\in\mathcal F_{m,\epsilon}(b)\). The defining moment equalities for \(\mathcal F_{m,\epsilon}(b)\), together with linearity of integration over the finite sum, give
\[
\int h\,d\nu
=\sum_{c\in\mathcal I_m}\gamma_c\int B_c^{(m)}\,d\nu
=\sum_{c\in\mathcal I_m}\gamma_c b_c.
\]
The last expression depends only on \(b\), so \(\int h\,d\nu\) is identical for all measures in every feasible fiber.

Conversely, suppose that \(h\notin V_m\). We first exhibit a point in the relative interior of \(\mathcal B_{m,\epsilon}\). The set \(K_{\epsilon}\) is nonempty: the four-cell vector \(q^{\circ}=(1/4,1/4,1/4,1/4)\) has \(p(q^{\circ})=1/2\), and hence belongs to \(K_{\epsilon}\) because \(0<\epsilon<1/2\). Let \(d\) be the dimension of the affine hull of \(\mathcal B_{m,\epsilon}\). Since
\[
\mathcal B_{m,\epsilon}=\operatorname{conv}\{g_m(q):q\in K_{\epsilon}\}
\]
by the degree-\(m\) Bernstein moment-body definition, one may choose \(q_0,\ldots,q_d\in K_{\epsilon}\) such that \(g_m(q_0),\ldots,g_m(q_d)\) are affinely independent and span that affine hull. The barycenter
\[
b_0=\frac{1}{d+1}\sum_{k=0}^{d}g_m(q_k)
\]
is feasible. To verify directly that it is relatively interior, use the affine coordinates determined by this affine basis: all \(d+1\) barycentric coordinates of \(b_0\) equal \(1/(d+1)>0\), so a sufficiently small relative neighborhood of \(b_0\) keeps every coordinate positive and is contained in the simplex \(\operatorname{conv}\{g_m(q_0),\ldots,g_m(q_d)\}\), which in turn is contained in \(\mathcal B_{m,\epsilon}\). Thus \(b_0\in\operatorname{relint}(\mathcal B_{m,\epsilon})\).

By \(\mathrm{lem{:}full\text{-}support\text{-}perturbation}\), there is \(\nu_0\in\mathcal F_{m,\epsilon}(b_0)\) whose support is all of \(K_{\epsilon}\). Regard \(V_m\) as a subspace of \(L^2(\nu_0)\). Its image is finite dimensional and therefore closed. Moreover, if \(v\in V_m\) is zero \(\nu_0\)-almost surely, then continuity of \(v\) and full support of \(\nu_0\) imply \(v=0\) everywhere on \(K_{\epsilon}\). Consequently, if \(h\) agreed \(\nu_0\)-almost surely with a member of \(V_m\), continuity and full support would imply that it agreed with that member everywhere, contrary to \(h\notin V_m\).

Let \(P_0h\) be the \(L^2(\nu_0)\)-orthogonal projection of \(h\) onto \(V_m\), and set \(r=h-P_0h\). Both terms are continuous on the compact set \(K_{\epsilon}\), so \(r\) is bounded and continuous. Orthogonality and the definition of \(V_m\) yield
\[
\int B_c^{(m)}r\,d\nu_0=0
\qquad(c\in\mathcal I_m).
\]
The multinomial identity gives
\[
\sum_{c\in\mathcal I_m}B_c^{(m)}(q)
=\left(\sum_{a,y}q_{ay}\right)^m=1
\qquad(q\in K_{\epsilon}),
\]
so the constant function \(1\) belongs to \(V_m\), and therefore \(\int r\,d\nu_0=0\). Furthermore,
\[
\int hr\,d\nu_0
=\int (P_0h+r)r\,d\nu_0
=\int r^2\,d\nu_0>0.
\]
The strict inequality follows because equality would imply \(r=0\) \(\nu_0\)-almost surely, which by the preceding full-support argument would imply \(h\in V_m\).

Since \(r\ne0\), its supremum norm is positive. Choose \(s\) with \(0<s<\|r\|_{\infty}^{-1}\), and define two Borel probability measures on \(K_{\epsilon}\) by
\[
d\nu_0^{\pm}=(1\pm sr)\,d\nu_0.
\]
Indeed, \(1\pm sr>0\) everywhere by the choice of \(s\), and
\[
\int(1\pm sr)\,d\nu_0=1
\]
because \(\int r\,d\nu_0=0\). For every \(c\in\mathcal I_m\),
\[
\int B_c^{(m)}\,d\nu_0^{\pm}
=\int B_c^{(m)}\,d\nu_0
\pm s\int B_c^{(m)}r\,d\nu_0
=b_{0,c}.
\]
Hence \(\nu_0^+,\nu_0^-\in\mathcal F_{m,\epsilon}(b_0)\). On the other hand,
\[
\int h\,d\nu_0^+-\int h\,d\nu_0^-
=2s\int hr\,d\nu_0
=2s\int r^2\,d\nu_0>0.
\]
Thus \(\int h\,d\nu\) cannot be identical throughout every feasible fiber when \(h\notin V_m\). This proves the equivalence.

For \(h=f\), \(\mathrm{lem{:}finite\text{-}plate\text{-}nonpolynomial}\) gives \(f\notin V_m\) for every finite \(m\ge1\), so the equivalence implies failure of universal point identification. Finally, let \(b\in\operatorname{relint}(\mathcal B_{m,\epsilon})\). The final assertion of \(\mathrm{lem{:}full\text{-}support\text{-}perturbation}\), applied to \(f\notin V_m\), supplies \(\nu_b^+,\nu_b^-\in\mathcal F_{m,\epsilon}(b)\) with
\[
\int f\,d\nu_b^-\ne\int f\,d\nu_b^+.
\]
By the definitions of \(L_m(b)\) and \(U_m(b)\) as the infimum and supremum over this fiber,
\[
L_m(b)\le
\min\left\{\int f\,d\nu_b^-,\int f\,d\nu_b^+\right\}
<
\max\left\{\int f\,d\nu_b^-,\int f\,d\nu_b^+\right\}
\le U_m(b).
\]
Therefore \(L_m(b)<U_m(b)\) at every relative-interior feasible count law, as claimed.